\documentclass[11pt,a4paper]{article}
\usepackage[utf8]{inputenc}
\usepackage[T1]{fontenc}
\usepackage[spanish,es-minimal,es-noshorthands]{babel}
\usepackage{lmodern}
\usepackage[sfdefault,lining]{FiraSans}
\usepackage[a4paper,top=2.3cm,bottom=2.2cm,left=1.8cm,right=1.8cm,headheight=1.3cm,headsep=0.4cm,footskip=1cm]{geometry}
\usepackage{amsmath,amssymb}
\usepackage{icomma}
\usepackage{graphicx}
\usepackage[table]{xcolor}
\usepackage{tikz}
\usetikzlibrary{calc}
\usepackage{fontawesome5}
\usepackage{ccicons}
\usepackage{fancyhdr}
\usepackage{lastpage}
\usepackage{tasks}
\usepackage{needspace}
\usepackage[most]{tcolorbox}
\usepackage{array}
\usepackage[hidelinks]{hyperref}
\AtBeginDocument{\hypersetup{pdftitle={Ficha de ejercicios de la unidad 3: Sistemas de ecuaciones e inecuaciones. 1º Bach.},pdfsubject={Matemáticas I 1º Bach.},pdfkeywords={sistemas de ecuaciones, Gauss, sistemas no lineales, inecuaciones, sistemas de inecuaciones, regiones del plano}}}

\definecolor{navy}{HTML}{1B2748}
\definecolor{accent}{HTML}{E8603A}
\definecolor{teal}{HTML}{1F9E90}
\definecolor{soft}{HTML}{EEF1F8}
\definecolor{muted}{HTML}{5D6685}

\setlength{\parindent}{0pt}
\setlength{\parskip}{3pt}
\settasks{label=\alph*), label-format=\bfseries\color{navy}, label-width=1.3em, item-indent=1.8em, column-sep=0.6em, after-item-skip=5pt, before-skip=3pt, after-skip=4pt}

% ---------- capçalera i peu ----------
\newcommand{\llicencia}{\href{https://creativecommons.org/licenses/by-nc-sa/4.0/deed.es}{\ccbyncsa}\ \ Licencia Creative Commons Atribución-NoComercial-CompartirIgual 4.0 Internacional (CC BY-NC-SA 4.0)}
\pagestyle{fancy}
\fancyhf{}
\fancyhead[L]{\raisebox{-0.35\height}{\includegraphics[height=0.95cm]{logo-gva.png}}}
\fancyhead[C]{\small\color{muted}Ficha de ejercicios · Unidad 3: Sistemas de ecuaciones e inecuaciones · 1\textsuperscript{o} Bach.}
\fancyhead[R]{\raisebox{-0.35\height}{\includegraphics[height=0.95cm]{logo-vermellar.png}}}
\fancyfoot[L]{\scriptsize\color{muted}\llicencia}
\fancyfoot[R]{\small\color{muted}\thepage}
\renewcommand{\headrulewidth}{0.4pt}
\renewcommand{\headrule}{\color{muted!40}\hrule height \headrulewidth}
\fancypagestyle{primera}{\fancyhf{}\renewcommand{\headrulewidth}{0pt}
  \fancyfoot[L]{\scriptsize\color{muted}\llicencia}
  \fancyfoot[R]{\small\color{muted}\thepage}}

% ---------- nivells i exercicis ----------
\newcommand{\estrelles}[1]{{\color{accent}\ifcase#1\or\faStar\faStar[regular]\faStar[regular]\or\faStar\faStar\faStar[regular]\or\faStar\faStar\faStar\fi}}
\newcounter{exer}
\newcommand{\ex}[1]{\par\needspace{4\baselineskip}\medskip\refstepcounter{exer}\noindent{\color{navy}\bfseries\theexer.}\ \ }

\newcommand{\seccio}[2]{\par\needspace{16\baselineskip}\bigskip
  \begin{tcolorbox}[colback=navy,colframe=navy,coltext=white,arc=2mm,boxrule=0pt,left=3mm,right=3mm,top=1.2mm,bottom=1.2mm]
  {\large\bfseries #1\quad #2}
  \end{tcolorbox}\vspace{-2pt}}

\newcommand{\exf}[1]{\par\medskip\refstepcounter{exer}\noindent{\color{navy}\bfseries\theexer.}\ \ }
\providecommand{\lt}{<}\providecommand{\gt}{>}
\newcolumntype{C}[1]{>{\centering\arraybackslash}p{#1}}
% quadre de recordatori (només a la fitxa d'exercicis)
\newcommand{\recorda}[1]{\par\smallskip\begin{tcolorbox}[colback=accent!7,colframe=accent!55,arc=2mm,boxrule=0.5pt,left=2.5mm,right=2.5mm,top=1mm,bottom=1mm]\small{\color{accent}\faLightbulb[regular]}\ \textbf{Recuerda:}\ #1\end{tcolorbox}}
\newcommand{\res}[1]{#1}\newcommand{\rest}[1]{#1}

\def\fmneed{11}% que el títol d'una secció no quede sol al final del full
% =====================================================================
% Capçalera FentMates de les fitxes: un color per a cada curs.
% Els colors són els de la pàgina inicial (LEVELS a assets/js/materials.js);
% el logo ve de _eines/logo/pdf/ (python3 _eines/logo/genera.py).
%
% Ús, al final del preàmbul de la unitat:
%   % =====================================================================
% Capçalera FentMates de les fitxes: un color per a cada curs.
% Els colors són els de la pàgina inicial (LEVELS a assets/js/materials.js);
% el logo ve de _eines/logo/pdf/ (python3 _eines/logo/genera.py).
%
% Ús, al final del preàmbul de la unitat:
%   % =====================================================================
% Capçalera FentMates de les fitxes: un color per a cada curs.
% Els colors són els de la pàgina inicial (LEVELS a assets/js/materials.js);
% el logo ve de _eines/logo/pdf/ (python3 _eines/logo/genera.py).
%
% Ús, al final del preàmbul de la unitat:
%   \input{../comu/fentmates.tex}
%   \fmcurs{1eso}
% i, al començament del document:
%   \capcalera{Matemàtiques}{1r ESO}{Fitxa d'exercicis · Unitat 2}{Divisibilitat}
% Opcional abans de \input: \def\fmneed{14} (espai mínim abans de cada secció).
% =====================================================================
\makeatletter
\def\fm@logodir{../../logo/pdf}
\providecommand{\fmneed}{10}
\def\fm@cursnom{}

% clau del curs, color i color del text damunt del color
\newcommand{\fm@curs}[3]{\@namedef{fm@color@#1}{#2}\@namedef{fm@text@#1}{#3}}
\fm@curs{1eso}{E6B84A}{navy}
\fm@curs{2eso}{3B82F6}{white}
\fm@curs{3eso}{FF7A50}{white}
\fm@curs{4eso-a}{A78BFA}{white}
\fm@curs{4eso-b}{6D28D9}{white}
\fm@curs{1batx-mi}{9B1C1C}{white}
\fm@curs{1batx-ccss}{EF6B6B}{white}
\fm@curs{2batx-mii}{1E6B47}{white}
\fm@curs{2batx-ccss}{2BB3A3}{white}

\newcommand{\fmcurs}[1]{%
  \@ifundefined{fm@color@#1}{\PackageError{fentmates}{Curs desconegut: #1}{}}{}%
  \edef\fm@tmp{\noexpand\definecolor{curs}{HTML}{\@nameuse{fm@color@#1}}%
               \noexpand\colorlet{curstext}{\@nameuse{fm@text@#1}}}\fm@tmp
  \colorlet{cursfosc}{curs!72!black}%
  \def\fm@logo{\fm@logodir/vertical-#1.pdf}%
  \def\fm@marca{\fm@logodir/marca-#1.pdf}}

% ---------- capçalera de la primera pàgina ----------
% #1 matèria (petit), #2 curs, #3 tipus de document i unitat, #4 títol
\newcommand{\capcalera}[4]{%
  \gdef\fm@cursnom{#2}%
  \par\noindent
  \begin{tcolorbox}[colback=curs,colframe=curs,boxrule=0pt,arc=2.2mm,width=3.5cm,height=2.15cm,
      valign=center,left=3mm,right=2mm,top=1mm,bottom=1mm,nobeforeafter,box align=center]
    \color{curstext}{\footnotesize #1}\par\vspace{1pt}{\fontsize{19}{21}\selectfont\bfseries #2\par}%
  \end{tcolorbox}\hspace{5mm}%
  \begin{minipage}[c]{\dimexpr\linewidth-3.5cm-5mm-2.9cm\relax}
    {\color{cursfosc}\small\bfseries\MakeUppercase{#3}}\par\vspace{2pt}
    {\color{navy}\fontsize{20}{23}\selectfont\bfseries #4\par}
  \end{minipage}\hfill
  \begin{minipage}[c]{2.6cm}\raggedleft\includegraphics[width=2.6cm]{\fm@logo}\end{minipage}\par
  \vspace{2pt}}

% ---------- barres de secció del color del curs ----------
\renewcommand{\seccio}[2]{\par\needspace{\fmneed\baselineskip}\bigskip
  \begin{tcolorbox}[colback=curs,colframe=curs,coltext=curstext,arc=2mm,boxrule=0pt,left=3mm,right=3mm,top=1.2mm,bottom=1.2mm]
  {\large\bfseries #1\quad #2}
  \end{tcolorbox}\vspace{-2pt}}

% ---------- estrelles de dificultat del color del curs (a Batxillerat no se'n posen) ----------
\@ifundefined{estrelles}{}{%
  \renewcommand{\estrelles}[1]{{\color{curs}\ifcase#1\or\faStar\faStar[regular]\faStar[regular]\or\faStar\faStar\faStar[regular]\or\faStar\faStar\faStar\fi}}}

% ---------- capçalera i peu de les pàgines ----------
% primera pàgina: logos institucionals al peu
\fancypagestyle{primera}{\fancyhf{}\renewcommand{\headrulewidth}{0pt}
  \fancyfoot[L]{\raisebox{-0.32\height}{\includegraphics[height=0.78cm]{logo-gva.png}}\hspace{5mm}%
    \raisebox{-0.32\height}{\includegraphics[height=0.78cm]{logo-vermellar.png}}}
  \fancyfoot[R]{\scriptsize\color{muted}\href{https://creativecommons.org/licenses/by-nc-sa/4.0/deed.ca}{\ccbyncsa}\ \ CC BY-NC-SA 4.0 · fentmates.org\qquad\small\thepage}}
% la resta: el símbol a l'esquerra i el curs a la dreta, amb el títol al centre com abans
\fancyhead[L]{\raisebox{-0.3\height}{\includegraphics[height=0.62cm]{\fm@marca}}}
\fancyhead[R]{\small\bfseries\color{cursfosc}\fm@cursnom}
\renewcommand{\headrule}{\color{curs!45}\hrule height \headrulewidth}
\makeatother

%   \fmcurs{1eso}
% i, al començament del document:
%   \capcalera{Matemàtiques}{1r ESO}{Fitxa d'exercicis · Unitat 2}{Divisibilitat}
% Opcional abans de \input: \def\fmneed{14} (espai mínim abans de cada secció).
% =====================================================================
\makeatletter
\def\fm@logodir{../../logo/pdf}
\providecommand{\fmneed}{10}
\def\fm@cursnom{}

% clau del curs, color i color del text damunt del color
\newcommand{\fm@curs}[3]{\@namedef{fm@color@#1}{#2}\@namedef{fm@text@#1}{#3}}
\fm@curs{1eso}{E6B84A}{navy}
\fm@curs{2eso}{3B82F6}{white}
\fm@curs{3eso}{FF7A50}{white}
\fm@curs{4eso-a}{A78BFA}{white}
\fm@curs{4eso-b}{6D28D9}{white}
\fm@curs{1batx-mi}{9B1C1C}{white}
\fm@curs{1batx-ccss}{EF6B6B}{white}
\fm@curs{2batx-mii}{1E6B47}{white}
\fm@curs{2batx-ccss}{2BB3A3}{white}

\newcommand{\fmcurs}[1]{%
  \@ifundefined{fm@color@#1}{\PackageError{fentmates}{Curs desconegut: #1}{}}{}%
  \edef\fm@tmp{\noexpand\definecolor{curs}{HTML}{\@nameuse{fm@color@#1}}%
               \noexpand\colorlet{curstext}{\@nameuse{fm@text@#1}}}\fm@tmp
  \colorlet{cursfosc}{curs!72!black}%
  \def\fm@logo{\fm@logodir/vertical-#1.pdf}%
  \def\fm@marca{\fm@logodir/marca-#1.pdf}}

% ---------- capçalera de la primera pàgina ----------
% #1 matèria (petit), #2 curs, #3 tipus de document i unitat, #4 títol
\newcommand{\capcalera}[4]{%
  \gdef\fm@cursnom{#2}%
  \par\noindent
  \begin{tcolorbox}[colback=curs,colframe=curs,boxrule=0pt,arc=2.2mm,width=3.5cm,height=2.15cm,
      valign=center,left=3mm,right=2mm,top=1mm,bottom=1mm,nobeforeafter,box align=center]
    \color{curstext}{\footnotesize #1}\par\vspace{1pt}{\fontsize{19}{21}\selectfont\bfseries #2\par}%
  \end{tcolorbox}\hspace{5mm}%
  \begin{minipage}[c]{\dimexpr\linewidth-3.5cm-5mm-2.9cm\relax}
    {\color{cursfosc}\small\bfseries\MakeUppercase{#3}}\par\vspace{2pt}
    {\color{navy}\fontsize{20}{23}\selectfont\bfseries #4\par}
  \end{minipage}\hfill
  \begin{minipage}[c]{2.6cm}\raggedleft\includegraphics[width=2.6cm]{\fm@logo}\end{minipage}\par
  \vspace{2pt}}

% ---------- barres de secció del color del curs ----------
\renewcommand{\seccio}[2]{\par\needspace{\fmneed\baselineskip}\bigskip
  \begin{tcolorbox}[colback=curs,colframe=curs,coltext=curstext,arc=2mm,boxrule=0pt,left=3mm,right=3mm,top=1.2mm,bottom=1.2mm]
  {\large\bfseries #1\quad #2}
  \end{tcolorbox}\vspace{-2pt}}

% ---------- estrelles de dificultat del color del curs (a Batxillerat no se'n posen) ----------
\@ifundefined{estrelles}{}{%
  \renewcommand{\estrelles}[1]{{\color{curs}\ifcase#1\or\faStar\faStar[regular]\faStar[regular]\or\faStar\faStar\faStar[regular]\or\faStar\faStar\faStar\fi}}}

% ---------- capçalera i peu de les pàgines ----------
% primera pàgina: logos institucionals al peu
\fancypagestyle{primera}{\fancyhf{}\renewcommand{\headrulewidth}{0pt}
  \fancyfoot[L]{\raisebox{-0.32\height}{\includegraphics[height=0.78cm]{logo-gva.png}}\hspace{5mm}%
    \raisebox{-0.32\height}{\includegraphics[height=0.78cm]{logo-vermellar.png}}}
  \fancyfoot[R]{\scriptsize\color{muted}\href{https://creativecommons.org/licenses/by-nc-sa/4.0/deed.ca}{\ccbyncsa}\ \ CC BY-NC-SA 4.0 · fentmates.org\qquad\small\thepage}}
% la resta: el símbol a l'esquerra i el curs a la dreta, amb el títol al centre com abans
\fancyhead[L]{\raisebox{-0.3\height}{\includegraphics[height=0.62cm]{\fm@marca}}}
\fancyhead[R]{\small\bfseries\color{cursfosc}\fm@cursnom}
\renewcommand{\headrule}{\color{curs!45}\hrule height \headrulewidth}
\makeatother

%   \fmcurs{1eso}
% i, al començament del document:
%   \capcalera{Matemàtiques}{1r ESO}{Fitxa d'exercicis · Unitat 2}{Divisibilitat}
% Opcional abans de \input: \def\fmneed{14} (espai mínim abans de cada secció).
% =====================================================================
\makeatletter
\def\fm@logodir{../../logo/pdf}
\providecommand{\fmneed}{10}
\def\fm@cursnom{}

% clau del curs, color i color del text damunt del color
\newcommand{\fm@curs}[3]{\@namedef{fm@color@#1}{#2}\@namedef{fm@text@#1}{#3}}
\fm@curs{1eso}{E6B84A}{navy}
\fm@curs{2eso}{3B82F6}{white}
\fm@curs{3eso}{FF7A50}{white}
\fm@curs{4eso-a}{A78BFA}{white}
\fm@curs{4eso-b}{6D28D9}{white}
\fm@curs{1batx-mi}{9B1C1C}{white}
\fm@curs{1batx-ccss}{EF6B6B}{white}
\fm@curs{2batx-mii}{1E6B47}{white}
\fm@curs{2batx-ccss}{2BB3A3}{white}

\newcommand{\fmcurs}[1]{%
  \@ifundefined{fm@color@#1}{\PackageError{fentmates}{Curs desconegut: #1}{}}{}%
  \edef\fm@tmp{\noexpand\definecolor{curs}{HTML}{\@nameuse{fm@color@#1}}%
               \noexpand\colorlet{curstext}{\@nameuse{fm@text@#1}}}\fm@tmp
  \colorlet{cursfosc}{curs!72!black}%
  \def\fm@logo{\fm@logodir/vertical-#1.pdf}%
  \def\fm@marca{\fm@logodir/marca-#1.pdf}}

% ---------- capçalera de la primera pàgina ----------
% #1 matèria (petit), #2 curs, #3 tipus de document i unitat, #4 títol
\newcommand{\capcalera}[4]{%
  \gdef\fm@cursnom{#2}%
  \par\noindent
  \begin{tcolorbox}[colback=curs,colframe=curs,boxrule=0pt,arc=2.2mm,width=3.5cm,height=2.15cm,
      valign=center,left=3mm,right=2mm,top=1mm,bottom=1mm,nobeforeafter,box align=center]
    \color{curstext}{\footnotesize #1}\par\vspace{1pt}{\fontsize{19}{21}\selectfont\bfseries #2\par}%
  \end{tcolorbox}\hspace{5mm}%
  \begin{minipage}[c]{\dimexpr\linewidth-3.5cm-5mm-2.9cm\relax}
    {\color{cursfosc}\small\bfseries\MakeUppercase{#3}}\par\vspace{2pt}
    {\color{navy}\fontsize{20}{23}\selectfont\bfseries #4\par}
  \end{minipage}\hfill
  \begin{minipage}[c]{2.6cm}\raggedleft\includegraphics[width=2.6cm]{\fm@logo}\end{minipage}\par
  \vspace{2pt}}

% ---------- barres de secció del color del curs ----------
\renewcommand{\seccio}[2]{\par\needspace{\fmneed\baselineskip}\bigskip
  \begin{tcolorbox}[colback=curs,colframe=curs,coltext=curstext,arc=2mm,boxrule=0pt,left=3mm,right=3mm,top=1.2mm,bottom=1.2mm]
  {\large\bfseries #1\quad #2}
  \end{tcolorbox}\vspace{-2pt}}

% ---------- estrelles de dificultat del color del curs (a Batxillerat no se'n posen) ----------
\@ifundefined{estrelles}{}{%
  \renewcommand{\estrelles}[1]{{\color{curs}\ifcase#1\or\faStar\faStar[regular]\faStar[regular]\or\faStar\faStar\faStar[regular]\or\faStar\faStar\faStar\fi}}}

% ---------- capçalera i peu de les pàgines ----------
% primera pàgina: logos institucionals al peu
\fancypagestyle{primera}{\fancyhf{}\renewcommand{\headrulewidth}{0pt}
  \fancyfoot[L]{\raisebox{-0.32\height}{\includegraphics[height=0.78cm]{logo-gva.png}}\hspace{5mm}%
    \raisebox{-0.32\height}{\includegraphics[height=0.78cm]{logo-vermellar.png}}}
  \fancyfoot[R]{\scriptsize\color{muted}\href{https://creativecommons.org/licenses/by-nc-sa/4.0/deed.ca}{\ccbyncsa}\ \ CC BY-NC-SA 4.0 · fentmates.org\qquad\small\thepage}}
% la resta: el símbol a l'esquerra i el curs a la dreta, amb el títol al centre com abans
\fancyhead[L]{\raisebox{-0.3\height}{\includegraphics[height=0.62cm]{\fm@marca}}}
\fancyhead[R]{\small\bfseries\color{cursfosc}\fm@cursnom}
\renewcommand{\headrule}{\color{curs!45}\hrule height \headrulewidth}
\makeatother

\fmcurs{1batx-mi}

\begin{document}
\thispagestyle{primera}

\capcalera{Matemáticas I}{1\textsuperscript{o} Bach.}{Ficha de ejercicios · Unidad 3}{\raggedright Sistemas de ecuaciones e inecuaciones}
\vspace{10pt}
Nombre y apellidos: \rule{7.8cm}{0.4pt}\hfill Grupo: \rule{1.6cm}{0.4pt}\hfill Fecha: \rule{2.2cm}{0.4pt}
\vspace{8pt}

\begin{tcolorbox}[colback=soft,colframe=soft,arc=2mm,left=3mm,right=3mm,top=1.5mm,bottom=1.5mm]
\small Resuelve los ejercicios en el cuaderno mostrando todos los pasos. Da los resultados de manera exacta y, si hace falta, también aproximados. En los sistemas con denominadores, radicales o logaritmos, comprueba siempre las soluciones. Los ejercicios 1--28 son de la ficha del departamento; la sección 6 tiene otros nuevos para ampliar y aplicar lo que has aprendido.
\end{tcolorbox}

% =====================================================
\seccio{1}{Sistemas de ecuaciones lineales. Método de Gauss}

\exf{1} Resuelve los siguientes sistemas de ecuaciones lineales:\begin{tasks}(2)
  \task $\left\{\begin{array}{l}y=\dfrac{x+1}{2}+3\\y=2x+10\end{array}\right.$
  \task $\left\{\begin{array}{l}2(2x+y)-3(3x-2y)=-34\\\dfrac x2-\dfrac y3=2\end{array}\right.$
\end{tasks}

\ex{2} Resuelve los siguientes sistemas de ecuaciones lineales por el método de Gauss y clasifícalos:\begin{tasks}(2)
  \task $\left\{\begin{array}{l}x+2y-3z=7\\2x+y-z=6\\3x-y-z=6\end{array}\right.$
  \task $\left\{\begin{array}{l}x+y-z=5\\3x-2y-2z=5\\2x-3y-z=3\end{array}\right.$
  \task $\left\{\begin{array}{l}2x+y+z=5\\2x-y+3z=3\\3x+2y+z=8\end{array}\right.$
  \task $\left\{\begin{array}{l}x+y-2z=0\\2x-3y+3z=4\\5x-5y+4z=8\end{array}\right.$
  \task $\left\{\begin{array}{l}5x+2y-2z=0\\3x-y+3z=0\\8x+y+z=-1\end{array}\right.$
  \task $\left\{\begin{array}{l}x+3y-2z=0\\x-y+3z=0\\2x+2y+z=0\end{array}\right.$
  \task $\left\{\begin{array}{l}2y-z=-1\\5x-y-3z=2\\x-y+2z=3\end{array}\right.$
  \task $\left\{\begin{array}{l}x+2y-2z=2\\3x-3y+z=-14\\5x-y-2z=-15\end{array}\right.$
  \task $\left\{\begin{array}{l}x+y+z=1\\2x-3z=5\\2y+5z=2\end{array}\right.$
  \task $\left\{\begin{array}{l}x+3y-2z+2w=12\\2x-2y-z+w=5\\3x+y-2z-4w=16\\3x-3z-3w=15\end{array}\right.$
\end{tasks}\recorda{Haz ceros debajo de la diagonal. La última fila $(0\ \ 0\ \ m\,|\,n)$ decide: $m\neq0$ SCD; $m=n=0$ SCI; $m=0,\ n\neq0$ SI.}


% =====================================================
\seccio{2}{Sistemas de ecuaciones no lineales}

\exf{2} Resuelve los siguientes sistemas de ecuaciones no lineales:\begin{tasks}(3)
  \task $\left\{\begin{array}{l}2x+3y=-4\\x^2-2xy=5\end{array}\right.$
  \task $\left\{\begin{array}{l}\frac1x+\frac2y=4\\10xy=9\end{array}\right.$
  \task $\left\{\begin{array}{l}x^2+2y^2=6\\3x^2-y^2=11\end{array}\right.$
  \task $\left\{\begin{array}{l}2x+y=8\\2x+3y^2=22\end{array}\right.$
  \task $\left\{\begin{array}{l}2x+5y=-8\\xy-3y=4\end{array}\right.$
  \task $\left\{\begin{array}{l}x+y=0\\xy=1\end{array}\right.$
  \task $\left\{\begin{array}{l}2x-y=3\\x^2-y^2=3\end{array}\right.$
  \task $\left\{\begin{array}{l}x+y=11\\\frac1x+\frac6{y+1}=1\end{array}\right.$
  \task $\left\{\begin{array}{l}xy=-3\\x^2+2y^2=19\end{array}\right.$
  \task $\left\{\begin{array}{l}2x^2+3y^2=32\\-3x^2+4y^2=-48\end{array}\right.$
  \task $\left\{\begin{array}{l}\sqrt x+\sqrt y=15\\x-y=105\end{array}\right.$
  \task $\left\{\begin{array}{l}x^2-2xy=16-y^2\\x^2-y^2=72\end{array}\right.$
\end{tasks}

\ex{2} Resuelve los siguientes sistemas de ecuaciones exponenciales:\begin{tasks}(2)
  \task $\left\{\begin{array}{l}2^x+5^y=9\\2^{x+2}+5^{y+1}=41\end{array}\right.$
  \task $\left\{\begin{array}{l}2^x+3^y=5\\2^{x+1}-3^y=7\end{array}\right.$
  \task $\left\{\begin{array}{l}2^x-2^y=1016\\2^{x-y}=128\end{array}\right.$
  \task $\left\{\begin{array}{l}3^x+7^y=16\\3^{x-1}-7^{y+2}=-340\end{array}\right.$
  \task $\left\{\begin{array}{l}\dfrac{1}{2^x}+\dfrac{3}{2^y}=\dfrac{x+y}{8}\\x+y=5\end{array}\right.$
\end{tasks}\recorda{Cambio de variable: $u=2^x$, $v=5^y$\ldots\ y resuelve el sistema lineal en $u$ y $v$.}

\ex{2} Halla las soluciones de los siguientes sistemas de ecuaciones logarítmicas:\begin{tasks}(2)
  \task $\left\{\begin{array}{l}2\log x-3\log y=7\\\log x+\log y=1\end{array}\right.$
  \task $\left\{\begin{array}{l}x-y=21\\\log x+\log y=2\end{array}\right.$
  \task $\left\{\begin{array}{l}\log x+3\log y=5\\\log x-\log y=3\end{array}\right.$
  \task $\left\{\begin{array}{l}\log_3x=7-\log_3y\\x-2y=27\end{array}\right.$
  \task $\left\{\begin{array}{l}3x+2y=64\\\log x-\log y=1\end{array}\right.$
  \task $\left\{\begin{array}{l}\log x^2-2\log y=6\\2\log x+\log y^2=2\end{array}\right.$
\end{tasks}\recorda{Con $u=\log x$, $v=\log y$ se obtiene un sistema lineal; o bien agrupa los logaritmos. Comprueba que $x\gt0$ e $y\gt0$.}


% =====================================================
\seccio{3}{Sistemas de inecuaciones con una incógnita}

\exf{2} Resuelve los sistemas de inecuaciones siguientes:\begin{tasks}(3)
  \task $\left\{\begin{array}{l}2x-3\le0\\3x+1\gt0\end{array}\right.$
  \task $\left\{\begin{array}{l}x+2\le2\\2x-\frac{x-1}{3}\le x-\frac13\\x^2+2x-8\gt0\end{array}\right.$
  \task $\left\{\begin{array}{l}2x+1\lt x+2\\3x-1\le4x\end{array}\right.$
  \task $\left\{\begin{array}{l}3x-4\gt x\\x\ge2x-1\end{array}\right.$
  \task $\left\{\begin{array}{l}x^2+2x-3\ge0\\2(x-3)\gt-6\end{array}\right.$
  \task $\left\{\begin{array}{l}\dfrac{x-1}{x+2}\gt0\\x^2\ge1\end{array}\right.$
\end{tasks}


% =====================================================
\seccio{4}{Inecuaciones y sistemas con dos incógnitas}

\exf{1} Resuelve las siguientes inecuaciones con dos incógnitas:\begin{tasks}(2)
  \task $3x+4y\le12$
  \task $-3x+4y\ge4$
  \task $2x+y\lt2$
  \task $x\gt y-2$
\end{tasks}\recorda{Dibuja la recta (continua si es $\le,\ge$; discontinua si es $\lt,\gt$) y prueba un punto que no esté en ella, como $O(0,0)$.}

\ex{2} Resuelve los sistemas de inecuaciones con dos incógnitas:\begin{tasks}(3)
  \task $\left\{\begin{array}{l}0\le y\lt2\\x+y\le3\\x-y\ge2\\x\ge0,\ y\ge0\end{array}\right.$
  \task $\left\{\begin{array}{l}y\lt-2x+4\\y\ge x\end{array}\right.$
  \task $\left\{\begin{array}{l}y\le2\\x\gt-4\\x-y\lt-3\end{array}\right.$
  \task $\left\{\begin{array}{l}x+y\le6\\2y\le x\\x\ge0\end{array}\right.$
  \task $\left\{\begin{array}{l}2x+y\lt2\\x\gt y-2\\x\gt0,\ y\gt0\end{array}\right.$
  \task $\left\{\begin{array}{l}3\le y\le5\\x-6\le0\\x-4\ge0\end{array}\right.$
\end{tasks}


% =====================================================
\seccio{5}{Problemas de sistemas e inecuaciones}

\exf{1} En una granja hay 250 animales entre conejos y gallinas. Si el número total de patas es 800, ¿cuántos animales hay de cada tipo?

\ex{3} Hemos necesitado 10 dm$^2$ de cartón para construir una caja con tapa, de base cuadrada y 2 dm$^3$ de volumen. ¿Cuáles son las dimensiones de la caja si el lado de la base es un número entero de decímetros?

\ex{2} En una caja registradora hay billetes de 50\,€, 100\,€ y 200\,€: 21 billetes en total y 1\,800\,€. El número de billetes de 50\,€ es el quíntuple de los de 200\,€. ¿Cuántos hay de cada clase?

\ex{2} Tres nietos quieren hacer un regalo de 60\,€ a su abuela: Luis, el mayor, aporta el triple de lo que aportan los otros dos juntos, y Carmen aporta 3\,€ por cada 2 que aporta Pedro. ¿Cuánto aporta cada uno?

\ex{2} Calcula cuántas bicicletas, triciclos y coches tiene una juguetería si en total hay 50 vehículos, el número de bicicletas y triciclos cuadruplica el de coches y hay 140 ruedas.

\ex{2} Para un espectáculo hay entradas A, B y C, de 10, 15 y 30\,€. La recaudación fue de 775\,€, se vendieron 5 entradas más del tipo B que del C y el 20\,\% de las entradas eran del tipo C. ¿Cuántas se vendieron de cada tipo?

\ex{2} Un cajero solo tiene billetes de 10, 20 y 50\,€. Sacamos 290\,€ y nos da exactamente 8 billetes; el número de billetes de 10\,€ es el doble de los de 20\,€. Plantea y resuelve el sistema.

\ex{2} Un cajero automático contiene 800 billetes de 10, 20 y 50\,€ con un importe de 21\,000\,€. El número de billetes de 10\,€ es igual que el de 20\,€ y 50\,€ juntos. ¿Cuántos billetes hay de cada tipo?

\ex{3} En un cine se han vendido 1405 entradas y se han recaudado 7920\,€. La entrada normal cuesta 6\,€, la del día del espectador 4\,€ y la de los jubilados 3\,€. La recaudación de las entradas de precio reducido es el 10\,\% de la de las entradas normales. ¿Cuántas entradas de cada tipo se han vendido?

\ex{2} Un comerciante vende relojes A a 200\,€, B a 500\,€ y C a 250\,€. Un mes vende 200; los del tipo B son tantos como los de A y C juntos, y la recaudación es de 73\,500\,€. ¿Cuántos ha vendido de cada tipo?

\ex{2} Una persona compra memorias externas, lectores de libros electrónicos y tabletas gráficas a 100, 120 y 150\,€ la unidad. Paga 1160\,€ por 9 unidades y compra una tableta más que lectores. ¿Cuántas unidades de cada producto compra?

\ex{2} Después de un descuento del 10\,\% en cada precio, se ha pagado 3,96\,€ por un rotulador, un cuaderno y una carpeta. El cuaderno vale la mitad que el rotulador y la carpeta, lo que vale el cuaderno más el 20\,\% del rotulador. Calcula el precio original de cada objeto.

\ex{3} Una persona invierte 7000\,€ en acciones de las empresas A y B y en un depósito al 1\,\%. Con las acciones obtiene un 5\,\% (A) y un 3\,\% (B), y el beneficio total es de 202\,€. El dinero en acciones supera en 2600\,€ al del depósito. ¿Cuánto destinó a cada inversión?

\ex{2} Una agencia de viajes vende 60 billetes a Londres, París y Roma. Los de París son el doble de los otros dos destinos juntos y a Roma se emiten dos más que la mitad de los de Londres. ¿Cuántos billetes se han vendido para cada destino?

\ex{3} Un comerciante compra 200 kg de melocotones, 100 de manzanas y 300 de peras. Los vende subiendo un 25\,\% el precio de los melocotones y las manzanas y un 40\,\% el de las peras, y obtiene 1087\,€, de los cuales 257 son de beneficio. Si el kilo de melocotones le costó 50 céntimos más que el de peras, ¿cuál fue el precio de compra de cada fruta?

\ex{3} Un restaurante ofrece desayunos a 4\,€, comidas a 8\,€ y cenas a 10\,€. El último sábado se sirvieron tantas comidas como desayunos y cenas juntos, la recaudación total fue de 1116\,€ y la de las comidas superó a la de las cenas en 156\,€.\begin{tasks}(1)
  \task ¿Cuántos desayunos, comidas y cenas se sirvieron?
  \task ¿Qué beneficio se obtuvo si las ganancias son de 2,5\,€ por desayuno, 4\,€ por comida y 5\,€ por cena?
\end{tasks}

\ex{2} Un estudiante sacó 7,5 puntos en un examen de tres preguntas. En la tercera sacó un punto más que en la segunda, y los puntos de la primera quintuplicaban la diferencia entre los de la tercera y la primera. ¿Qué puntuación sacó en cada pregunta?

\ex{3} Un inversor reparte 42\,000\,€ entre una cuenta de ahorro al 5\,\%, un depósito a plazo fijo al 7\,\% y unos bonos al 9\,\% anual. En un año obtiene 2600\,€ de intereses, y los de la cuenta de ahorro son 200\,€ menos que la suma de los de las otras dos inversiones. ¿Cuánto invirtió en cada producto?

\ex{2} El largo de un rectángulo es 1 cm más que el doble del ancho. Queremos que el perímetro sea como mínimo de 20 cm y como máximo de 30 cm. ¿Entre qué valores puede estar el ancho?

\ex{2} Un artesano fabrica collares y pulseras. Un collar necesita 2 horas de trabajo y 1 unidad de material, y una pulsera, 1 hora y 2 unidades de material. Cada semana dispone, como máximo, de 40 horas y 50 unidades de material.\begin{tasks}(1)
  \task Plantea el sistema de inecuaciones y representa la región de las posibles combinaciones de collares y pulseras.
  \task ¿Puede fabricar 15 collares y 12 pulseras en una semana? ¿Y 10 collares y 20 pulseras?
\end{tasks}


% =====================================================
\seccio{6}{Ampliación y problemas}

\exf{3} Clasifica el sistema según los valores de $a$ y resuélvelo cuando sea compatible: \ $\left\{\begin{array}{l}x+y+z=3\\x-y+z=1\\2x+az=4\end{array}\right.$

\ex{2} Un rectángulo tiene 10 cm de perímetro y 6 cm$^2$ de área. ¿Cuáles son sus dimensiones?

\ex{2} Hace 2 años la edad de un padre era el triple de la del hijo y dentro de 10 años será el doble. ¿Cuántos años tienen ahora?

\ex{3} Escribe un sistema de inecuaciones cuya solución sea el triángulo de vértices $(0,0)$, $(4,0)$ y $(0,2)$, con los bordes incluidos.

\vfill
\begin{center}\color{muted}\small\faCheckCircle[regular]\ \ ¡Ánimo! Comprueba tus resultados con las herramientas interactivas de los apuntes de la unidad.\end{center}

\end{document}
